\documentclass{THUThesis}


\sufesetup{
    title = {金融统计 \LaTeX 模板},
    author = {},
    id = {},
    school ={},
    discipline = {},
} 

\begin{document}
%%%%定义一级小标题(section)，二级小标题(subsection)，三级小标题(subsubsection)字体字号， 四级"（1）"还未定义
\titleformat

\maketitle


\newpage
% \tableofcontents
\newpage
\setcounter{page}{1}%开始正文第一页

\section{序言}
序言应说明本课题的意义、目的、主要研究内容、范围及应解决的问题。
\subsection{研究背景及意义}

\subsection{研究现状与评价}


\subsection{本文贡献与创新点}

\subsection{文章基本思路与结构}


\section{论文格式概述}

\subsection{正文之前}

\subsubsection{封页}

封页上的内容一律按照统一封面的样张式样打印，必须正 确无误。题目用二号黑体字，其他用四号宋体字。

\subsubsection{题目和摘要}

作业题目为三号黑体字，可以分为1或2行 居中打印。作业题目下空一项居中打印“摘要”二字(三号 黑体)，字间空一格，标点符号占一格。“摘要”二字下空一行打印 内容(四号宋体)。每段开头空二格，标点符号占一格。摘要内容后 下空一行打印“关键字”三字(四号黑体字)，其后为关键词(四号 宋体)。英文摘要题目全部采用Arial字体，可分成1-3行居中打印， 以下均用五号Arial字体。摘要以浓缩的形式概括研究课题的内容，中文摘要在 300汉字左右，英文摘要应与中文摘要基本相对应。关键词是表述作业主题内容信息的单词或 术语，关键词数量一般不超过6个。每一个关键词之间用逗号隔 开，最后一个关键词后不用标点符号。

\subsubsection{目录}
“目录”二字(三号黑体)，下空二行为章、节、小节及其开始页码。目录是作业各组成部分的小标题，文字应简 明扼要。目录按章节排列编写，标明页数，便于阅读。章节、小 节分别以一、(一)1.(1)等数字依次标出。要求标题层次清 晰。目录中的标题应与正文中的标题一致。

\subsubsection{标题}
每章标题以三号字黑体居中打印;“章”下空二行为“节”，以四 号黑体左起打印;“节”下空一行为“小节”，以五号黑体左起打印。换行打印 作业正文。作业题目应该简短、明确、有概括性;字数 要适当，一般不宜超过20个汉字。

\subsection{正文格式要求}
\subsubsection{正文}
采用五号宋体字打印。正文是对研究工作的详细表述，一般由标题、文字、 图、表格和公式等部分组成。
\subsubsection{图}
图题若采用中英文对照时，其英文字体为五号正体，中文字体为五号 楷体。引用图应在图题的左上角标出文献来源;图号按章顺序编号，如图3-1为 第三章第一图。如图含有几个不同部分应将分图号标注在分图的左上角，并在图 题下列出各部分内容。\par
一张图片大概这样，简单如图 \ref{fig:hanhan} 所示。
\begin{figure}[!ht]
  \centering
  \includegraphics[width=0.45\linewidth]{图片1.jpg}
  \caption{实验数据的图片}
  \label{fig:hanhan}
\end{figure} 
多张图片，网上也有很多方法。
% 简单如图 \ref{fig:hanhan1} 所示。
% \begin{figure}
%   \centering
%   \includegraphics[width=0.45\linewidth]{图片1.jpg}
%   \includegraphics[width=0.45\linewidth]{图片1.jpg}
%   \includegraphics[width=0.45\linewidth]{图片1.jpg}
%   \includegraphics[width=0.45\linewidth]{图片1.jpg}
%   \caption{实验数据的图片}
%   \label{fig:hanhan1}
% \end{figure} 
\subsubsection{表格}
表格按章顺序编号，如表3-1为第三章第一表。表应有标题，表内必须按规定的符号注明单位。
对表格来说并排表格既可以像图~\ref{tab:parallel1}、
表\ref{tab:parallel2} 使用小页环境，也可以如表\ref{tab:subtable} 使用子表格来做。R语言可以使用xtable程序包直接导出表格代码，复制粘贴一下就ok。

\begin{table}[htbp]
\noindent\begin{minipage}{0.5\textwidth}
\centering
\caption{第一个并排子表格}
\label{tab:parallel1}
\begin{tabular}{p{2cm}p{2cm}}
\toprule
111 & 222 \\\midrule
222 & 333 \\\bottomrule
\end{tabular}
\end{minipage}%
\begin{minipage}{0.5\textwidth}
\centering
\caption{第二个并排子表格}
\label{tab:parallel2}
\begin{tabular}{p{2cm}p{2cm}}
\toprule
111 & 222 \\\midrule
222 & 333 \\\bottomrule
\end{tabular}
\end{minipage}
\end{table}


\begin{table}[htbp]
\centering
\caption{并排子表格}
\label{tab:subtable}
\subcaptionbox{第一个子表格}
{
\begin{tabular}{c|c|c}
\hline
模型 & 得分1 &得分2 \\
\hline
a &111 & 222 \\
b &222 & 333 \\
\hline
\end{tabular}
}
\hskip2cm
\subcaptionbox{第二个子表格}
{
\begin{tabular}{c|c|c}
\hline
模型 & 得分1 &得分2 \\
\hline
a &111 & 222 \\
b &222 & 333 \\
\hline
\end{tabular}
}
\end{table}


\subsubsection{公式}
公式书写应在文中另起一行。并使用英文标点。
考虑模型\eqref{eq1}
\begin{equation}
y_i = \mu(t_i) + e_i, \quad i = 1,2,\dots n,
\label{eq1}
\end{equation}
如果需要换行
\begin{equation}
  \begin{aligned}
  y_i &= \mu(t_i) + e_i\\
      &= \mu(t_i) + e_i.\\
  \end{aligned}
\end{equation}

\subsubsection{参考文献}
不可缺少的组成部分，也是作者对他人知识
成果的承认和尊重。参考文献一律列于文末。连续出版物(期刊):序号 作者，题名刊名，出版年，期号，起止 日;
专著:序号 作者，书名，版本(第1版不标注)，出版地，出版年， 起止页码; 论文集:序号，作者，题名，主编，论文集名，出版地，出版年，起 止页码;
学位论文:序号 题名;【学位论文】(英文用【Dissertation】),
保存地点，保存单位，年份。\par
使用方法：在导出bib格式文献引用，将内容复制到main.bib，之后直接使用“cite”，如文献中指出\cite{guegan2005can},\citet{guegan2005can}和\citet{dahlhaus1997}指出，文献中指出\cite{guegan2005can,dahlhaus1997}等用法。


\section{实证分析}
%S&P 500 daily returns. 
\subsection{数据介绍}
介绍数据背景相关文献，表格图片计算结果以及结合实际的深刻分析。

% \begin{figure}
%   \centering
%   \includegraphics[width=0.45\textwidth]{data/IBM_LR.png}
%   \includegraphics[width=0.45\textwidth]{data/IBM_SLR.png}
%   \caption{左图：IBM股票从2005年7月14日到2013年8月30日的对数收益率；右图：IBM股票从2005年7月14日到2013年8月30日的平方对数收益率}
% \end{figure}

\subsection{数据分析}


\section{结论与局限性}
\subsection{结论}
总结全文
\subsection{局限性}
不足和未来发展方向

\newpage

\bibstyle{gbt7714-numerical}
\bibliography{main}
 
\section{附录}
一些不想放在正文的细节。不宜放在正文中，但有参考价值的内容。如调查问卷、 公式推演、编写程序、原始数据附表等。
\subsection{核心代码}
\subsubsection{检验统计量计算}
可以直接从文件中导入代码。见注释。
% \lstinputlisting[language=r]{Statistic.R}



\end{document}
